\section{Related Works} \label{related_works}
% 将相关工作分成两个或多个部分分别阐述。每个部分都可以再按范畴来进一步细分。与上一章介绍的内容相比，这里对每组方法的优势与不足的总结要更加详细到位，文献引用的数量也应适当增加。而在逻辑上，作者仍然需要想好论文的中心思想是什么，综述内容一定要切题，从各个侧面来突显论文的贡献

% 研究现状一
\subsection{Short Text Clustering Algorithms based on Topic Models}
Topic models assume data points are generated by a mixture model, and then employ techniques like EM~\cite{bishop2006pattern} or Gibbs sampling~\cite{robert2004monte} to estimate the parameters of the mixture model, so as to obtain the clustering results.

Nigam et al.~\cite{nigam2000text} utilize the DMM model for classification with both labeled and unlabeled documents and propose an EM-based algorithm~\cite{nigam2000text} for DMM (abbr. to EM-DMM). When only unlabeled documents are provided, DMM turns out to be a clustering model. Yu et al.~\cite{YuDPM:KDD2010} propose the DPMFS model for text clustering that can also infer the number of clusters automatically. Due to the slow convergence of DPMFS~\cite{DBLP:journals/tkde/HuangYWZS13}, they further propose the DMAFP model as an approximation to the DPMFS and introduce a variational inference algorithm. They compare DMAFP with other four clustering models: EM-DMM~\cite{nigam2000text}, K-means~\cite{jain2010K-means}, LDA~\cite{Blei:LDA2003}, and EDCM~\cite{DCMelkan2006clustering}. Rangrej et al.~\cite{Rangrej:WWW2011} compare three clustering algorithms for short text clustering: K-means, singular value decomposition and affinity propagation on a small set of tweets, and find that Affinity Propagation~\cite{frey2007clustering} can achieve better performance than the other two algorithms. However, the complexity of Affinity Propagation is quadratic in the number of documents, which means Affinity Propagation cannot scale to huge datasets with millions of documents.
Banerjee et al.~\cite{Banerjee:SIGIR2007} propose a method of improving the accuracy of short text clustering by enriching their representation with additional features from Wikipedia. Topic models like LDA~\cite{Blei:LDA2003} and PLSA~\cite{hofmann1999probabilistic} are probabilistic generative models that can model texts and identify latent semantics underlying the text collection. 
LDA has excellent performance on normal-length texts, but it is inferior on short texts because LDA depends on cross-text word co-occurrences, which are lacking in short texts. 
% 方法的优势与不足的总结，我们方法如何解决
Therefore, we propose GSDMM to address the short text clustering problem and manage the high dimensionality of short texts, achieving excellent performance in clustering short texts. Like Topic Models~\cite{Blei:LDA2003}, GSDMM can also obtain the representative words of each cluster. This enables GSDMM to efficiently address key clustering challenges, including the sparsity issue in short texts. Furthermore, previous methods overlooked the importance of individual words to text clustering, thereby oversimplifying the nature of textual data. In reality, words exhibit varying levels of discriminative power across categories. Therefore, \sysname{} is proposed to adaptively adjust the influence of individual words based on their informational significance within each cluster. This approach highlights discriminative words while minimizing the impact of less informative terms, effectively capturing subtle local patterns.

% Some models are proposed to remedy this issue. BTM~\cite{BTM} introduces the concept of biterm into topic model, which is an unordered word pair in a short text and models the generation process of biterms instead of short texts. SATM~\cite{SATM} views the generation process of short texts as two phases. Stand topic models are employed to generate regular-sized texts from which short texts are sampled. PTM~\cite{PTM} considers there exist latent pseudo texts composed of short texts and models the topic distribution on long pseudo-texts instead of short texts. Zuo et al~\cite{WNTM} proposed to create a word co-occurrences network and infers the topic distribution of each word instead of each short text.

% With the development of word embedding technology, researchers have sought to incorporate pre-trained word embeddings into topic models to alleviate the sparsity of short texts. Firstly, the LF-DMM~\cite{nguyen2015improving} model introduces pre-trained word embeddings into GSDMM, effectively improving clustering performance. Later, Li et al.~\cite{li2016topic} introduced a GPU strategy in GSDMM that modifies the collapsed Gibbs Sampling process. The GPU strategy entails that when a word $w$ is added to topic $z$, similar words to $w$ are also added to topic $z$. The similarity scores of words are obtained through pre-trained word embedding computations. Recently, Li et al.~\cite{li2017enhancing} proposed the GPU-PDMM model, which applies the GPU strategy to the PDMM model. The PDMM model uses the Poisson distribution to modify the assumption of GSDMM from "one short text is generated by only one topic" to "one short text can be generated by a limited number of topics." 

% 研究现状二
\subsection{Deep Short Text Clustering Algorithms}
Text clustering has been an enduring research focus. Many early deep text clustering algorithms employ various autoencoder to learn deep representations for clustering. Methods in the former category perform clustering in a two-stage manner, which first feeds texts into a deep neural network to learn the text representations, and then a certain clustering method, like K-means, is performed based on the learned representations. Xu et al~\cite{stcc} propose a self-taught learning framework constituted of a convolutional neural network. Guan et al.~\cite{guan2020deep} use the neural language model and InferSent~\cite{conneau2017supervised} to obtain short text embeddings and then apply K-means for clustering. The two-stage deep clustering methods separate the optimization process into the representations learning step and clustering step, where the objective of representation learning is not consistent with clustering sometimes, leading to unsatisfactory results.

Methods in the latter category perform representation learning and text clustering in an end-to-end manner, thereby allowing joint optimization of the two steps. ARL~\cite{arl} integrates these two components into a unified model that constructs low-dimensional text representations from one-hot encoding through word embedding and mean pooling. Subsequently, attention techniques~\cite{attention} are leveraged between text representations and cluster representations. DECCA~\cite{DECCA} learns text representations using a contrastive autoencoder, upon which a deep embedding clustering framework is built. Guan et al.~\cite{guan2020deep} propose a deep feature-based text clustering model named DFTC, which incorporates text encoders and an explanation module into text clustering tasks. SCCL~\cite{SCCL} develops a novel contrastive learning framework that pulls together the learned representations of positive pairs augmented from the same text while pushing apart negative pairs from other texts. SCCL utilizes the Sentence-BERT~\cite{conneau2017supervised} model and fine-tunes it using contrastive loss and clustering loss to obtain high-quality text embeddings. As is well known, efficiency is extremely important in clustering tasks. 
% 方法的优势与不足的总结，我们方法如何解决
However, past methods, whether introducing embedding representations or various deep models, have incurred significant efficiency losses while improving effectiveness. Building upon GSDMM and \sysname{}, we have ensured high efficiency while greatly enhancing the performance of short text clustering.
